\subsection{Corpora, Genre and Collocations}
\label{sec:corpora}
\slidesrc{Terminology introduction: 39--41}

\subsubsection{Corpus}
\slidesrc{Terminology introduction: 39}

\begin{nbdef}{Corpus (plural: corpora)}
A \textbf{text corpus} is a large and structured set of texts (nowadays usually stored
and processed electronically). Corpora are used for statistical analysis and hypothesis
testing, for checking occurrences, and for validating linguistic rules within a specific
language territory.
\end{nbdef}

Example: the corpora by Mark Davies, originally at \url{https://corpus.byu.edu/} and now
at \url{https://www.english-corpora.org/}: COCA (Corpus of Contemporary American English),
NOW (News on the Web), and others. You can search them to see how often and in which
contexts a word or phrase is used.

\begin{nbwhy}
Corpora are the raw material of statistical and neural NLP
(Section~\ref{sec:nlp-history}): probabilities, word embeddings and LLMs are all learnt
from corpora. What the corpus contains (its size, languages, genres, time period)
determines what the model learns.
\end{nbwhy}

\subsubsection{Genre}
\slidesrc{Terminology introduction: 40}

\begin{nbdef}{Genre}
\begin{itemize}[nosep]
  \item \textbf{Genre}: a style or category of art, music or literature.
  \item \textbf{Genre of a corpus}: usually defined on the basis of
        \emph{text-external} features, such as medium, function and format (e.g.\
        newspaper articles, academic papers, fiction, spoken conversation transcripts,
        social media posts).
\end{itemize}
\end{nbdef}

``Text-external'' means that the genre is decided by \emph{where the text comes from and
what it is for}, not by counting words inside it. Texts of different genres differ a
lot (sentence length, vocabulary, formality), so a model trained on newspapers often
performs worse on tweets.

\subsubsection{Collocations}
\slidesrc{Terminology introduction: 41}

\begin{nbdef}{Collocation}
A \textbf{collocation} is a sequence of words or terms that co-occur \emph{more often
than would be expected by chance}. Examples: ``pay attention'', ``fast food'', ``make an
effort'', ``powerful engine''.
\end{nbdef}

\emph{Fast} and \emph{quick} mean roughly the same thing, yet native speakers combine
them with different nouns (example from the NOW corpus):

\begin{center}
\begin{tabular}{@{}l l@{}}
\toprule
\textbf{Natural English} & \textbf{Unnatural English} \\
\midrule
the fast train  & the \sout{quick} train \\
fast food       & \sout{quick} food \\
a quick shower  & a \sout{fast} shower \\
a quick meal    & a \sout{fast} meal \\
\bottomrule
\end{tabular}
\end{center}

\begin{nbintuition}[``more often than by chance'']
If words were combined at random, \emph{fast food} and \emph{quick food} would appear
about as often as the frequencies of \emph{fast} and \emph{quick} predict. In a corpus,
\emph{fast food} appears far more often than that, and \emph{quick food} hardly ever.
That excess is what makes it a collocation. Nothing in grammar or meaning forbids
``quick food'': collocations are a matter of \emph{usage}, which is why they are learnt
from corpora.
\end{nbintuition}
